% !TEX program = xelatex
% 金华十校2026年11月高三模拟考试 语文参考答案与评分标准
\documentclass[11pt,a4paper]{ctexart}
\usepackage{ywexam}

\examfooter{语文参考答案}

\begin{document}

\answertitle{金华十校2026年11月高三模拟考试}{参考答案与评分标准}

\sechead{一、阅读（71分）}

\subsechead{（一）阅读Ⅰ（19分）}
\ans{1}{B。用《太玄经》这类后世抽象概念反推上古思维，属于原文所批判的“时代错置”误区。}
\ans{2}{C。“主要因为”解读有误，原文中“神话概念被投射到高山”还有一个重要原因，“高山因其高耸入云、最接近太阳（飞鸟）的视觉特征。”}
\ans{3}{A。大汶口陶尊的“日—火—山”符号，直接展现了上古先民对太阳与山地关联的具象观察，能够证明“昆仑观念源于对日鸟运行的观察”观点，其余选项分别与神话、地貌想象、金石有关，无法直接证明该观点。}
\ans{4}{\xq{1}首先提出“昆仑”是中华文化核心精神符号的主张；\xq{2}接着从反面驳斥昆仑文化研究的错误方式；\xq{3}然后正面从先民思维特点切入，完成“昆仑文化”的演进分析；\xq{4}最后总结“昆仑文化”三层世代层累完整演进过程。}
\scorenote{（每点1分，意思相近即可）}
\ans{5}{\xq{1}“太阳驿站”数字模拟活动；\xq{2}将河姆渡“双鸟朝阳”象牙蝶形器等文物以虚拟动画形式向参观者展示；\xq{3}让青少年直观感受昆仑文化从神话想象到落地现实的过程。}
\scorenote{（每点2分，意思相近即可）}

\subsechead{（二）阅读Ⅱ（16分）}
\ans{6}{C。C项“她已把补锅匠视为精神知己，并在他身上看到了丈夫所缺少的浪漫气质”属于过度推断。伊丽莎眼睛发亮，主要是因为自己的劳动成果被外界需要，并由此唤起她对更广阔世界的向往；她羡慕的是补锅匠所代表的漂泊生活，并未把他个人视为“精神知己”，文本中也没有将她丈夫与“浪漫气质”作对比的依据。}
\ans{7}{D。D项“已下定决心离开丈夫和牧场”属于无中生有。句\xq{4}是伊丽莎低声呢喃，表现她对远方、自由、广阔世界的向往和憧憬，但并未表明她已经下定决心出走。}
\ans{8}{\xq{1}以神态、动作等细节变化呈现心理转变。修锅人问起菊花，伊丽莎“眼睛闪闪发亮”，生出热烈期盼；结尾“无力地瘫在座位上”哭泣，走向怅惘失落。}
\subans{\xq{2}}{以菊花这个物象暗示心理变化。菊花是她对远方美好的寄托，满怀期待托付菊苗；后来看到菊苗被弃，物象的遭遇暗示期盼落空。}
\subans{\xq{3}}{以人物语言呈现心理转变。从“远方满是光明”到“我们要喝很多酒”呈现了从憧憬到失落的转变。}
\scorenote{（答出任意两点就得4分，其他答案只要言之有据就酌情给分。仅罗列“细节描写”等术语，缺少对心理转变过程的分析，最多得1分；必须结合前后情节，点明心理由热烈期盼走向失落，方可拿全分。）}
\ans{9}{观点一：赞扬女性的觉醒。}
\subans{理由：}{}
\subans{\xq{1}}{展示自我价值：伊丽莎能种出大菊花，确认自身能力与价值；穿着盛装追问丈夫“美”“健壮”的含义，展现女性自我意识觉醒。}
\subans{\xq{2}}{向往自由生活：伊丽莎主动向补锅匠倾诉内心憧憬，不甘被局限于家庭主妇身份，说“真希望女人也可以过这种生活”，在枯燥的家庭生活中，借栽种菊花萌生对美、远方与自由的向往。}
\subans{观点二：哀叹女性命运的不幸。}{}
\subans{理由：}{}
\subans{\xq{1}}{男性世界漠视女性精神追求：丈夫亨利只夸伊丽莎“健壮”，不懂她对美与自我的追求；补锅匠也直言旅途不属于女性，她的精神诉求不被时代接纳。}
\subans{\xq{2}}{女性觉醒后依然没有出路：伊丽莎寄托全部希望的菊苗被丢弃，幻想破碎，结尾她借酒麻醉，瘫坐遮掩“像老妇人似的虚弱地哭泣”，表明她无力挣脱社会与家庭对女性的束缚，命运充满无奈与悲哀。}
\scorenote{（答出一点3分，两点6分，其他答案只要言之有理就酌情给分）}

\subsechead{（三）阅读Ⅲ（21分）}
\ans{10}{ADG（敌者，所伐之国也，后虽无复，何伤哉？）}
\ans{11}{D。两词含义分别为“追击”，“使……跟随”，意思与用法均不同。}
\ans{12}{B。孔子评价晋文公“既知一时之权，又知万世之利”，说明其并未否定舅犯意见。}
\ans{13}{（1）用诡诈之术对待百姓，（纵使能）出其不意获取一时胜利，以后肯定不会有第二次。}
\scorenote{（“遇”“偷”“复”每词解释各1分，句意通顺1分。）}
\subans{（2）}{我们理亏而楚国理直，楚国的士兵没有不（为此）产生愤激情绪的，不能说他们已经士气衰颓。}
\scorenote{（“曲”“生气”“老”每词解释各1分，句意通顺1分。）}
\ans{14}{（1）《韩非子》：文公采用舅犯诈谋。《国语》：文公听从子犯劝谏，重诺讲礼，占据道义制高点（或“楚国得寸进尺，失礼理屈”）。（2）法家：重功利，尚权术（或“视克敌为国之大利”）。儒家：重德义，讲礼仪。}
\scorenote{（每一小点各1分。）}

\subsechead{（四）阅读Ⅳ（9分）}
\ans{15}{A。诗歌情感复杂深沉，“喜”字非诗眼，全诗的情感基调亦非“喜”。}
\ans{16}{读fèn。理由：\xq{1}尾联意为作者与卢纶素有情谊，两人的表亲关系加深了这种感情，“分”字解为“情分、情谊”，所以读fèn；\xq{2}fèn为仄声，与下句末字的平声构成呼应，符合五律平仄要求，音韵更和谐。}
\scorenote{（答出正确读音给1分；答出理由\xq{1}给3分，其中解说尾联句意2分，解说“分”字词义1分；答出理由\xq{2}给2分，只要说明音韵和谐，均可给分。如答读为fēn，且解说能自圆其说，酌情给1—3分；仅答读fēn，未加解说，不给分。）}

\subsechead{（五）名篇名句默写（6分）}
\ans{17}{（1）落木千山天远大，澄江一道月分明。}
\subans{（2）}{移船相近邀相见，添酒回灯重开宴。}
\subans{（3）}{示例：巴山楚水凄凉地，二十三年弃置身。（刘禹锡《酬乐天扬州初逢席上见赠》）}
\subans{}{塞上长城空自许，镜中衰鬓已先斑。（陆游《书愤》）}
\subans{}{欲渡黄河冰塞川，将登太行雪满山。（李白《行路难》）}
\scorenote{（每空1分，共6分。）}

\sechead{二、语言文字运用（19分）}
\ans{18}{“描摩”改为“描摹”；“须要”改为“需要”。}
\scorenote{（每改对一处得1分，未改对不得分。）}
\ans{19}{C。\xq{3}句表现“失意”时抄写作品的直接感受；\xq{4}句“三遍”承接\xq{5}句的“继续”；\xq{2}句的“淡”承接\xq{1}句的“痛”。}
\ans{20}{\xq{1}将范仲淹的豁达心境比作“暖流”，让抽象的状态描述成具体可感的体验。}
\subans{\xq{2}}{用“渗”“堵”“松动”等词，使“豁达”“郁结”等抽象概念有动态感。}
\subans{\xq{3}}{用“团”这个量词来表现“郁结”，使无形的情绪形态化。}
\scorenote{（每点2分，意思相近即可，其他角度酌情赋分）}
\ans{21}{（1）语句\xq{2}修改为：会让每个字慢下来。（2）语句\xq{6}修改为：只有笔尖与纸面的摩擦声。}
\ans{22}{\xq{1}释义：实现精神突围的手抄经典行为。}
\subans{依据：}{当代人通过抄书来疏解情绪、获得精神力量，体现传统抄写在当代的延伸。}
\subans{\xq{2}}{释义：收获独处空间的慢书写方式。}
\subans{依据：}{在碎片化信息的时代，抄书是人们脱离喧嚣、获得内心安静的新方式。}
\scorenote{（每点3分，其中，释义2分，依据1分，意思相近即可，其他角度酌情赋分）}

\sechead{三、写作（60分）}
\ans{23}{审题参考：材料共有三句话，围绕“激情”一词分层表述，存在内在的逻辑关系。第一句阐述了激情的价值；第二句指出年青人对于激情存在不同状况；第三句话提出年青人应该学会如何正确对待激情。}
\subans{立意参考：}{年青人要激情满怀地去创造价值，始终充满激情；但要张弛有度，不过度透支生命、不失理智冷静。}

\par\vspace{0.5em}
\sechead{文言文参考译文}

\translabel{材料一：}
\begin{translation}
晋文公将要和楚国人作战，召见舅犯问他，说：“我将要和楚国人作战，他们人多我们人少，对此该怎么办？”舅犯说：“我听说：讲究礼仪的君子，不嫌多施忠信；作战布阵之间，不嫌多施诡诈伪饰。您还是用诡诈的办法对付他们吧。”文公辞别舅犯，于是召见雍季问他，说：“我将要和楚国人作战，他们人多我们人少，对此该怎么办？”雍季回答说：“用诡诈之术对待百姓，（纵使能）出其不意获取一时胜利，以后肯定不会有第二次。”文公说：“好。”辞别雍季，用舅犯的计谋和楚国人作战，打败了他们。回来后行赏，先赏雍季而后赏舅犯。群臣说：“城濮之战，（采用的）是舅犯的计谋。采用他的言论，行赏时却把他排在后面，可以吗？”文公说：“这不是你们所能知道的。舅犯的话，是一时的权宜之计；雍季的话，则符合国家长久的利益。”孔子听说这件事，说：“文公称霸，是应该的啊！他既知道一时的权宜之计，又知道国家长久的利益。”

有人说：雍季的回答，不切合文公的问题。凡是回答提问的人，要根据所问问题的大小缓急来回答。所问的问题高大，却用卑下狭隘的话回答，那么英明的君主是不会接受的。如今文公问“怎样以少敌众”，（雍季）却回答说“以后肯定不会有第二次”，这不是用来应答的内容。况且文公既不懂得一时的权宜之计，又不懂得国家长久的利益。作战胜利，那么国家安定、自身稳固，军队强大、威望树立，即使以后有第二次（出现相同的情况），也没有不会比这次胜利获益更大了，怎么还要担心长久的利益不来？作战不胜，那么国家灭亡、军队衰弱，自身死亡、名声消亡，免除今天的死亡都来不及，哪里有空等待长久的利益？

等待万世的利益，在于今天的胜利；今天的胜利，在于对敌人使诈；对敌人使诈，就是万世的利益罢了。所以说：雍季的回答，不切合文公的问题。并且文公又不懂得舅犯的话。舅犯所说的“不厌诈伪”，不是说要欺诈自己的百姓，而是说要欺诈自己的敌人。敌人，是要讨伐的国家，以后即使没有第二次，又有什么损害呢？舅犯说“讲究礼仪的君子，不厌弃忠信”，忠是用来爱护自己的下属的，信是用来不欺骗自己的百姓的。已经谈到爱护下属、不欺骗百姓了，还有比这更好言论吗？然而一定要说“出于诈伪”，是因为这是军事斗争的计谋。舅犯前面有善言，后面有战胜之功，所以舅犯有两件功劳却排在后面论赏，雍季没有一件功劳却先受赏，（这是错误的。）“文公称霸，不也是应该的吗？”孔子不懂得赏赐的道理啊。
\end{translation}

\translabel{材料二：}
\begin{translation}
子玉解除对宋国的包围，追击晋国军队。楚国已经摆好阵势，晋国军队后退三十里。军吏请示说：“国君躲避臣子，这是耻辱。况且楚军已经士气衰颓，必定会失败。为什么要后退？”子犯说：“你们忘了在楚国时的事吗？我听说：作战时，理直就士气昂扬，理亏就士气衰颓。还没有报答楚国的恩惠却去救援宋国，（这是）我们理亏而楚国理直，楚国的士兵没有不（为此）产生愤激情绪的，不能说他们已经士气衰颓。如果我们以国君的身份躲避臣子，而他们还不撤离，那他们也就理亏了。”于是后退九十里以躲避楚军。楚军军士想要停下来，子玉不肯，到了城濮，果然交战，楚军大败。君子说：“子犯善于用德义来劝勉。”
\end{translation}

\end{document}
